\clearpage
\begin{explanationbox}
\textbf{\textcolor{CExplanation}{一、为什么会出现“区间最值”问题？}}
\begin{center}\begin{tikzpicture}[font=\small,>=stealth,line join=round]
\draw[gray!40,rounded corners=3pt] (0,0) rectangle (14.2,8.3);\node at (7.1,7.700000000000001) {顶点虽低，却不在允许的区间中};
\begin{scope}[shift={(1.8777777777777778,2.2249999999999996)},x=1.111111111111111cm,y=0.575cm]\draw[->,gray] (-0.7,0)--(4.7,0) node[right]{$x$};\draw[->,gray] (0,-1)--(0,7) node[above]{$y$};\draw[gray,thick,solid,domain=-0.2:4.2,samples=180] plot(\x,{(\x-2)^2+1});\draw[blue,thick,very thick,domain=3:4,samples=180] plot(\x,{(\x-2)^2+1});\fill[orange!85!black] (3,2) circle(2.2pt);\fill[orange!85!black] (4,5) circle(2.2pt);\draw[violet,dashed] (2,1.05)--(2,6.5);\draw[orange!85!black,very thick] (3,-.5)--(4,-.5);\node[below] at (2,1){$(2,1)$};\node[right] at (3,2){$(3,2)$};\node[above left=3pt] at (4,5){$(4,5)$};\foreach \v in {2,3,4}{\node[below=5pt] at (\v,-.5){$\v$};}\end{scope}\node[anchor=west,align=left,text width=5.0cm] at (8.4,4.0) {$f=(x-2)^2+1$\\[9pt]允许 $x\in[3,4]$\\[9pt]$2\notin[3,4]$\\[9pt]最近 $x=3$：最小值2\\[9pt]最远 $x=4$：最大值5};
\node[text width=12.6cm,align=center] at (7.1,.65) {灰线仅供定位；蓝色粗线是允许段，不能到灰色顶点处取值。};
\end{tikzpicture}\end{center}



学过二次函数后，我们知道 $y=(x-2)^2+1$ 的最低点在 $x=2$，最低值为 $1$。但如果题目规定 $3\le x\le4$，你就\textbf{不能再选 $x=2$}。这时必须只看抛物线上 $3\le x\le4$ 的那一小段。

这好比在学校走廊里规定了“只能在第 3 米到第 4 米之间站立”：即使离教室门口最近的位置在第 2 米，你也不能站过去。\textbf{先看允许站在哪里，再谈最近和最远。}

\medskip

\end{explanationbox}
\clearpage
\begin{explanationbox}
\textbf{\textcolor{CExplanation}{二、先用一个初中生也能算出来的例子观察}}
\begin{center}\begin{tikzpicture}[font=\small,>=stealth,line join=round]
\draw[gray!40,rounded corners=3pt] (0,0) rectangle (14.2,8.5);\node at (7.1,7.9) {同一条曲线，取不同的区间，最值位置就会改变};
\node at (2.4,6.9) {$[0,4]$};\begin{scope}[shift={(1.1407407407407406,3.0500000000000003)},x=0.6296296296296295cm,y=0.45cm]\draw[->,gray] (-0.7,0)--(4.7,0) node[right]{$x$};\draw[->,gray] (0,-1)--(0,7) node[above]{$y$};\draw[gray,thick,solid,domain=-0.2:4.2,samples=160] plot(\x,{(\x-2)^2+1});\draw[blue,thick,very thick,domain=0:4,samples=160] plot(\x,{(\x-2)^2+1});\fill[orange!85!black] (2,1) circle(2pt);\fill[orange!85!black] (0,5) circle(2pt);\fill[orange!85!black] (4,5) circle(2pt);\end{scope}\node[align=center,text width=4cm] at (2.4,1.7) {最小1，在2\\最大5，在0与4};\node at (7.050000000000001,6.9) {$[3,4]$};\begin{scope}[shift={(5.790740740740741,3.0500000000000003)},x=0.6296296296296295cm,y=0.45cm]\draw[->,gray] (-0.7,0)--(4.7,0) node[right]{$x$};\draw[->,gray] (0,-1)--(0,7) node[above]{$y$};\draw[gray,thick,solid,domain=-0.2:4.2,samples=160] plot(\x,{(\x-2)^2+1});\draw[blue,thick,very thick,domain=3:4,samples=160] plot(\x,{(\x-2)^2+1});\fill[orange!85!black] (3,2) circle(2pt);\fill[orange!85!black] (4,5) circle(2pt);\end{scope}\node[align=center,text width=4cm] at (7.050000000000001,1.7) {最小2，在3\\最大5，在4};\node at (11.7,6.9) {$[0,1]$};\begin{scope}[shift={(10.44074074074074,3.0500000000000003)},x=0.6296296296296295cm,y=0.45cm]\draw[->,gray] (-0.7,0)--(4.7,0) node[right]{$x$};\draw[->,gray] (0,-1)--(0,7) node[above]{$y$};\draw[gray,thick,solid,domain=-0.2:4.2,samples=160] plot(\x,{(\x-2)^2+1});\draw[blue,thick,very thick,domain=0:1,samples=160] plot(\x,{(\x-2)^2+1});\fill[orange!85!black] (1,2) circle(2pt);\fill[orange!85!black] (0,5) circle(2pt);\end{scope}\node[align=center,text width=4cm] at (11.7,1.7) {最小2，在1\\最大5，在0};
\node[text width=12.6cm,align=center] at (7.1,.65) {三图同尺度；蓝色粗线才是各题允许取到的部分。};
\end{tikzpicture}\end{center}



取 $f(x)=(x-2)^2+1$。它的函数值由两件事组成：\textbf{离 $x=2$ 有多远，和这个距离的平方}。
\[
\begin{array}{c|ccccc}
x&0&1&2&3&4\\ \hline
|x-2|&2&1&0&1&2\\
f(x)&5&2&1&2&5
\end{array}
\]
\begin{itemize}[leftmargin=*,itemsep=0.3em]
\item 允许范围是 $[0,4]$：可以到 $x=2$，最低值是 $1$；两端都离它最远，最高值是 $5$。
\item 改成 $[3,4]$：最靠近 $2$ 的位置是 $3$，最远的位置是 $4$，所以最低值 $f(3)=2$，最高值 $f(4)=5$。
\item 改成 $[0,1]$：最靠近 $2$ 的位置是 $1$，最远的位置是 $0$，所以最低值 $f(1)=2$，最高值 $f(0)=5$。
\end{itemize}

\textbf{眼前一亮的地方：} 不同的区间，虽然是\textbf{同一条抛物线}，最高点、最低点却可能换位置；因为我们截取的是不同的一段图像。

\medskip

\end{explanationbox}
\clearpage
\begin{explanationbox}
\textbf{\textcolor{CExplanation}{三、图像：为什么不能想当然地选择顶点？}}
\begin{center}\begin{tikzpicture}[font=\small,>=stealth,line join=round]
\draw[gray!40,rounded corners=3pt] (0,0) rectangle (14.2,8.3);\node at (7.1,7.700000000000001) {顶点虽低，却不在允许的区间中};
\begin{scope}[shift={(1.8777777777777778,2.2249999999999996)},x=1.111111111111111cm,y=0.575cm]\draw[->,gray] (-0.7,0)--(4.7,0) node[right]{$x$};\draw[->,gray] (0,-1)--(0,7) node[above]{$y$};\draw[gray,thick,solid,domain=-0.2:4.2,samples=180] plot(\x,{(\x-2)^2+1});\draw[blue,thick,very thick,domain=3:4,samples=180] plot(\x,{(\x-2)^2+1});\fill[orange!85!black] (3,2) circle(2.2pt);\fill[orange!85!black] (4,5) circle(2.2pt);\draw[violet,dashed] (2,1.05)--(2,6.5);\draw[orange!85!black,very thick] (3,-.5)--(4,-.5);\node[below] at (2,1){$(2,1)$};\node[right] at (3,2){$(3,2)$};\node[above left=3pt] at (4,5){$(4,5)$};\foreach \v in {2,3,4}{\node[below=5pt] at (\v,-.5){$\v$};}\end{scope}\node[anchor=west,align=left,text width=5.0cm] at (8.4,4.0) {$f=(x-2)^2+1$\\[9pt]允许 $x\in[3,4]$\\[9pt]$2\notin[3,4]$\\[9pt]最近 $x=3$：最小值2\\[9pt]最远 $x=4$：最大值5};
\node[text width=12.6cm,align=center] at (7.1,.65) {灰线仅供定位；蓝色粗线是允许段，不能到灰色顶点处取值。};
\end{tikzpicture}\end{center}




\textbf{读图：} 灰色是完整曲线的定位参考，蓝色粗线是允许段，橙色粗线标出横坐标范围 $[3,4]$。顶点虽然低，却不在允许范围中；这道题只能在蓝色粗线对应的那一段图像上找答案。

\medskip

\end{explanationbox}
\clearpage
\begin{explanationbox}
\textbf{\textcolor{CExplanation}{四、为什么“离对称轴的距离”能决定函数值？}}
\begin{center}\begin{tikzpicture}[font=\small,>=stealth,line join=round]
\draw[gray!40,rounded corners=3pt] (0,0) rectangle (14.2,8.8);\node at (7.1,8.200000000000001) {到轴的距离相同，a的正负改变高低关系};
\node at (3.1500000000000004,7.3) {$f=(x-1)^2+1$};\begin{scope}[shift={(1.3125,3.190909090909091)},x=1.225cm,y=0.34545454545454546cm]\draw[->,gray] (-0.5,0)--(3.5,0) node[right]{$x$};\draw[->,gray] (0,-2)--(0,9) node[above]{$y$};\draw[blue,thick,solid,domain=0:3,samples=180] plot(\x,{(\x-1)^2+1});\fill[orange!85!black] (0,2) circle(2pt);\fill[orange!85!black] (1,1) circle(2pt);\fill[orange!85!black] (3,5) circle(2pt);\end{scope}\node[align=center,text width=5.8cm] at (3.1500000000000004,1.7) {$[0,3]$：近处最小1\\远端$x=3$，最大5};\node at (9.95,7.3) {$g=-(x-1)^2+4$};\begin{scope}[shift={(8.1125,3.190909090909091)},x=1.225cm,y=0.34545454545454546cm]\draw[->,gray] (-0.5,0)--(3.5,0) node[right]{$x$};\draw[->,gray] (0,-2)--(0,9) node[above]{$y$};\draw[teal,thick,solid,domain=0:3,samples=180] plot(\x,{-(\x-1)^2+4});\fill[orange!85!black] (0,3) circle(2pt);\fill[orange!85!black] (1,4) circle(2pt);\fill[orange!85!black] (3,0) circle(2pt);\end{scope}\node[align=center,text width=5.8cm] at (9.95,1.7) {$[0,3]$：近处最大4\\远端$x=3$，最小0};
\node[text width=12.6cm,align=center] at (7.1,.65) {两图同尺度、同轴$x=1$：左图越远越高，右图越远越低。};
\end{tikzpicture}\end{center}



把一般二次函数写成
\[
 f(x)=a(x-h)^2+k,\qquad h=-\frac{b}{2a},\quad k=f(h).
\]
其中 $(x-h)^2$ 是“$x$ 离对称轴有多远”的平方。对于同一条抛物线，$a$ 和 $k$ 不变，所以比较函数值时只需比较 $|x-h|$。

\begin{itemize}[leftmargin=*,itemsep=0.3em]
\item $a>0$：距离平方越大，$a(x-h)^2$ 越大，因此函数值越大。
\item $a<0$：距离平方越大，$a(x-h)^2$ 越小，因此函数值越小。
\end{itemize}

\textbf{注意：}“越远越大”只对开口向上的抛物线成立；开口向下时恰好相反。

\medskip

\end{explanationbox}
\clearpage
\begin{explanationbox}
\textbf{\textcolor{CExplanation}{五、为什么最远处一定是端点？}}
\begin{center}\begin{tikzpicture}[font=\small,>=stealth,line join=round]
\draw[gray!40,rounded corners=3pt] (0,0) rectangle (14.2,8.4);\node at (7.1,7.800000000000001) {证明一：区间不动，先在数轴上找最近与最远};
\begin{scope}[shift={(2.5,6.1)},x=1.15cm]\draw[->,gray] (-1.6,0)--(5.8,0);\draw[blue,line width=2pt] (0,0)--(4,0);\foreach \v in {0,4}{\fill[blue] (\v,0) circle(2pt);\node[below=5pt] at (\v,0){$\v$};}\fill[violet] (-1,0) circle(2.5pt);\node[above=2pt] at (-1,0){$h=-1$};\draw[orange!85!black,line width=2pt] (0,-.16)--(0,.16);\node[anchor=west] at (-1.6,.95){轴在左侧：最近$0$；最远4};\end{scope}\begin{scope}[shift={(2.5,4.1)},x=1.15cm]\draw[->,gray] (-1.6,0)--(5.8,0);\draw[blue,line width=2pt] (0,0)--(4,0);\foreach \v in {0,4}{\fill[blue] (\v,0) circle(2pt);\node[below=5pt] at (\v,0){$\v$};}\fill[violet] (2,0) circle(2.5pt);\node[above=2pt] at (2,0){$h=2$};\draw[orange!85!black,line width=2pt] (2,-.16)--(2,.16);\node[anchor=west] at (-1.6,.95){轴在区间内：最近$2$；最远0与4};\end{scope}\begin{scope}[shift={(2.5,2.1)},x=1.15cm]\draw[->,gray] (-1.6,0)--(5.8,0);\draw[blue,line width=2pt] (0,0)--(4,0);\foreach \v in {0,4}{\fill[blue] (\v,0) circle(2pt);\node[below=5pt] at (\v,0){$\v$};}\fill[violet] (5,0) circle(2.5pt);\node[above=2pt] at (5,0){$h=5$};\draw[orange!85!black,line width=2pt] (4,-.16)--(4,.16);\node[anchor=west] at (-1.6,.95){轴在右侧：最近$4$；最远0};\end{scope}
\node[text width=12.6cm,align=center] at (7.1,.65) {蓝段为[0,4]，橙色短标记是最近位置；轴在区间内时可以直接取轴的位置。};
\end{tikzpicture}\end{center}



在一段直直的数轴上，任意 $x\in[m,n]$ 都夹在 $m,n$ 之间。固定一个位置 $h$，你从区间内任意位置出发，要走得离 $h$ 最远，不可能比左右两端中较远的那个还远。因此\textbf{最远位置只用比较 $m,n$}，根本不用检查中间所有点。

但最近处不同：如果 $h$ 正好在区间里，最近处就是 $h$ 自身；如果不在，就只能退到离它近的那个端点。

\medskip

\end{explanationbox}
\clearpage
\begin{explanationbox}
\textbf{\textcolor{CExplanation}{六、特别好用的“区间中点”技巧}}
\begin{center}\begin{tikzpicture}[font=\small,>=stealth,line join=round]
\draw[gray!40,rounded corners=3pt] (0,0) rectangle (14.2,8.6);\node at (7.1,8.0) {比远近：看对称轴在区间中点的哪一侧};
\begin{scope}[shift={(2,6.3)},x=1.1cm]\draw[->,gray] (-.5,0)--(8.8,0);\draw[blue,thick] (0,0)--(8,0);\foreach \v in {0,8}{\fill[blue] (\v,0) circle(2pt);\node[below=4pt] at (\v,0){$\v$};}\draw[violet,dashed] (4,-.2)--(4,.15);\fill[orange!85!black] (3,0) circle(2.5pt);\node[above=6pt] at (3,0){$h=3$};\node[anchor=west] at (0,1){左端距离3，右端距离5：右端更远};\end{scope}\begin{scope}[shift={(2,4.4)},x=1.1cm]\draw[->,gray] (-.5,0)--(8.8,0);\draw[blue,thick] (0,0)--(8,0);\foreach \v in {0,8}{\fill[blue] (\v,0) circle(2pt);\node[below=4pt] at (\v,0){$\v$};}\draw[violet,dashed] (4,-.2)--(4,.15);\fill[orange!85!black] (6,0) circle(2.5pt);\node[above=6pt] at (6,0){$h=6$};\node[anchor=west] at (0,1){左端距离6，右端距离2：左端更远};\end{scope}\begin{scope}[shift={(2,2.5)},x=1.1cm]\draw[->,gray] (-.5,0)--(8.8,0);\draw[blue,thick] (0,0)--(8,0);\foreach \v in {0,8}{\fill[blue] (\v,0) circle(2pt);\node[below=4pt] at (\v,0){$\v$};}\draw[violet,dashed] (4,-.2)--(4,.15);\fill[orange!85!black] (4,0) circle(2.5pt);\node[above=6pt] at (4,0){$h=4$};\node[anchor=west] at (0,1){左端距离4，右端距离4：两端同远};\end{scope}
\node[text width=12.6cm,align=center] at (7.1,.65) {同尺度；紫线为中点4，轴位于中点时两端才同远。};
\end{tikzpicture}\end{center}



假设区间是 $[0,8]$，中点是 $4$：
\begin{itemize}[leftmargin=*,itemsep=0.25em]
\item 对称轴在 $x=3$：它偏向左端，右端 $8$ 一定更远（距离分别为 $3$ 和 $5$）。
\item 对称轴在 $x=6$：它偏向右端，左端 $0$ 一定更远（距离分别为 $6$ 和 $2$）。
\item 对称轴在 $x=4$：两端同样远，所以 $f(0)=f(8)$，无论抛物线开口向上还是向下。
\end{itemize}
\textbf{记法：轴偏左，右端远；轴偏右，左端远；轴在正中，两端同高。}

\medskip

\end{explanationbox}
\clearpage
\begin{explanationbox}
\textbf{\textcolor{CExplanation}{七、做题时到底先干什么？}}
\begin{center}\begin{tikzpicture}[font=\small,>=stealth,line join=round]
\draw[gray!40,rounded corners=3pt] (0,0) rectangle (14.2,8.4);\node at (7.1,7.800000000000001) {证明一：区间不动，先在数轴上找最近与最远};
\begin{scope}[shift={(2.5,6.1)},x=1.15cm]\draw[->,gray] (-1.6,0)--(5.8,0);\draw[blue,line width=2pt] (0,0)--(4,0);\foreach \v in {0,4}{\fill[blue] (\v,0) circle(2pt);\node[below=5pt] at (\v,0){$\v$};}\fill[violet] (-1,0) circle(2.5pt);\node[above=2pt] at (-1,0){$h=-1$};\draw[orange!85!black,line width=2pt] (0,-.16)--(0,.16);\node[anchor=west] at (-1.6,.95){轴在左侧：最近$0$；最远4};\end{scope}\begin{scope}[shift={(2.5,4.1)},x=1.15cm]\draw[->,gray] (-1.6,0)--(5.8,0);\draw[blue,line width=2pt] (0,0)--(4,0);\foreach \v in {0,4}{\fill[blue] (\v,0) circle(2pt);\node[below=5pt] at (\v,0){$\v$};}\fill[violet] (2,0) circle(2.5pt);\node[above=2pt] at (2,0){$h=2$};\draw[orange!85!black,line width=2pt] (2,-.16)--(2,.16);\node[anchor=west] at (-1.6,.95){轴在区间内：最近$2$；最远0与4};\end{scope}\begin{scope}[shift={(2.5,2.1)},x=1.15cm]\draw[->,gray] (-1.6,0)--(5.8,0);\draw[blue,line width=2pt] (0,0)--(4,0);\foreach \v in {0,4}{\fill[blue] (\v,0) circle(2pt);\node[below=5pt] at (\v,0){$\v$};}\fill[violet] (5,0) circle(2.5pt);\node[above=2pt] at (5,0){$h=5$};\draw[orange!85!black,line width=2pt] (4,-.16)--(4,.16);\node[anchor=west] at (-1.6,.95){轴在右侧：最近$4$；最远0};\end{scope}
\node[text width=12.6cm,align=center] at (7.1,.65) {蓝段为[0,4]，橙色短标记是最近位置；轴在区间内时可以直接取轴的位置。};
\end{tikzpicture}\end{center}



先画一条数轴，标出 $m,n,h$；接着看开口向上还是向下；最后找最近处和最远处，并代入求值。先判断位置，往往比直接算三个函数值更快。
\end{explanationbox}
